% Recovered verbatim from chapters_11_15_with_part/chapters_11_15_with_part.tex
% Canonical chapter 12
\chapter{Cosmology with Chromatic Sectors}

\section{Chromatic Energy Densities}

The chromatic fields contribute separate energy densities
\[
  \rho_\chi =
  \frac{1}{2} G^{ab} \dot{\chi}_a \dot{\chi}_b + U(\bm{\chi}),
\]
with pressures determined by the kinetic and potential balance.
Different chromatic modes can redshift at different rates, leading to distinct cosmological eras dominated by different bands.

\section{Chromatic Era Transitions}

Transitions between chromatic-dominated eras correspond to changes in which $\chi_a$ carries the dominant fraction of $\rho_\chi$.
Such transitions can imprint features in the expansion history and may leave signatures in primordial perturbations.

\section{Approach to the White-State}

As the universe contracts or energy becomes highly compressed, the chromatic amplitude vector $\bm{\chi}$ tends toward the symmetric configuration defining the white-state core.
This behaviour is central to nonsingular cosmological cycles in MK--43 Ultra.

